% =====================================================================
%  CHAPTER 1A: 蛋白质 — 参考：0310营养参考_extracted.txt (PPT原文)
% =====================================================================
\chapter{蛋白质}
\chapterintro{掌握蛋白质的生理作用与氮平衡（零/负/正）；掌握必需氨基酸（9种）、条件必需氨基酸和限制氨基酸；掌握氨基酸模式与蛋白质互补三原则；掌握蛋白质消化率、BV、NPU、PER、AAS/PDCAAS（难点*）；掌握人体蛋白质营养状况评价指标。}

\section{蛋白质的生理作用与氮平衡}

\begin{keybox}
\textbf{蛋白质的生理作用：}
\begin{enumerate}[leftmargin=*]
    \item 人体内细胞的主要成分：构成人体各组织器官（皮肤、肌肉、骨骼、内脏等）；补偿新陈代谢的消耗；修补组织损失
    \item 人体内重要物质的主要原料：激素、酶、抗体、血红蛋白
    \item 提供热量：每克蛋白质在体内可产生16.72千焦耳的能量，由蛋白质提供的能量占每日人需总能量的\imp{10\%—15\%}
\end{enumerate}

\textbf{氮平衡（Nitrogen Balance）：}反映机体摄入氮和排出氮的关系。
$$B = I - (U + F + S)$$
$B$：氮平衡；$I$：摄入氮；$U$：尿氮；$F$：粪氮；$S$：皮肤等氮损失
\begin{itemize}[leftmargin=*]
    \item \imp{零氮平衡}：正常成人
    \item \imp{负氮平衡}：老年人、饥饿过度节食、消耗性疾病
    \item \imp{正氮平衡}：生长发育期、孕期、恢复期、肌肉增长
\end{itemize}
\end{keybox}

\section{必需氨基酸与蛋白质分类}

\begin{keybox}
\textbf{必需氨基酸（Essential Amino Acid, EAA）：}人体不能合成或合成速度不能满足机体需要，必须从食物中获得的一类氨基酸。共有\imp{9种}：亮氨酸、异亮氨酸、赖氨酸、蛋氨酸、苯丙氨酸、苏氨酸、色氨酸、缬氨酸、组氨酸。其中\imp{组氨酸是婴儿的必需氨基酸}。

\textbf{条件必需氨基酸（Conditionally Essential AA）/ 半必需氨基酸：}半胱氨酸和酪氨酸在体内分别由蛋氨酸和苯丙氨酸转变而成。如果膳食中能直接提供这两种氨基酸，则人体对蛋氨酸和苯丙氨酸的需要可\imp{分别减少30\%和50\%}。代谢路径：蛋氨酸→同型半胱氨酸→胱硫醚→半胱氨酸。

\textbf{氨基酸模式（Pattern）：}蛋白质中各种必需氨基酸的构成比例。计算方法是将该种蛋白质中的色氨酸含量定为1，分别计算出其他必需氨基酸的相应比值。

\textbf{蛋白质分类：}完全蛋白质（优质蛋白，植物仅大豆）；半完全蛋白质；不完全蛋白质（胶原蛋白）。

\textbf{限制氨基酸（Limiting Amino Acid）：}当食物蛋白质中一种或几种必需氨基酸相对含量较低或缺乏，限制了食物蛋白质中的其它必需氨基酸被机体利用的程度，使其营养价值降低。含量最低的称为\imp{第一限制氨基酸}。\textbf{谷类缺赖氨酸，豆类缺蛋氨酸。}

\textbf{蛋白质互补三原则：}(1)食物生物学种属愈远愈好；(2)搭配种类愈多愈好；(3)食用时间愈近愈好，最好同时食用。
\end{keybox}

\section{蛋白质营养学评价（难点*）}

\begin{keybox}
\textbf{1. 蛋白质含量：}各种蛋白质含氮量稳定，约占蛋白质重量的\imp{16\%}。\imp{蛋白质含量 = 氮含量 × 6.25}。检测方法：凯氏定氮法（Kjedahl法）、双缩脲法、Folin-酚试剂法（Lowry法）、紫外吸收法、考马斯亮蓝法（Bradford法）。

\textbf{2. 蛋白质消化率：}消化率(\%) = 氮吸收量/摄入氮量 × 100
\begin{itemize}[leftmargin=*]
    \item \imp{真消化率(\%) = [食物氮 $-$ (粪氮 $-$ 粪代谢氮)] / 食物氮 × 100}
    \item \imp{表观消化率(\%) = (食物氮 $-$ 粪氮) / 食物氮 × 100}
\end{itemize}

\textbf{3. 生物价（BV）：}反映食物蛋白质消化吸收后被机体利用程度的指标。BV越高，利用程度越高。
$$\text{BV} = \frac{\text{储留氮}}{\text{吸收氮}} \times 100,\quad \text{储留氮} = \text{吸收氮} - (\text{尿氮} - \text{尿内源性氮})$$

\textbf{4. 蛋白质净利用率（NPU）：}
$$\text{NPU}(\%) = \text{消化率} \times \text{生物价} = \frac{\text{储留氮量}}{\text{食物氮}} \times 100$$

\textbf{5. 蛋白质功效比值（PER）：}测定生长发育中的幼小动物摄入1g蛋白质所增加的体重克数。一般选用刚断乳的雄性大鼠，被测蛋白质作为唯一蛋白质来源占饲料10\%，喂饲28天。PER = 动物体重增加(g)/摄入蛋白质(g)。校正PER = 待测PER/酪蛋白PER × 2.5。

\textbf{6. 氨基酸评分（AAS）/ 蛋白质化学评分：}将待评食物蛋白质必需氨基酸与人体必需氨基酸模式比较。
$$\text{AAS} = \frac{\text{待测蛋白质每克氮中某种氨基酸含量(mg)}}{\text{理想模式每克氮中该氨基酸含量(mg)}} \times 100$$
$$\text{经消化率修正的AAS (PDCAAS)} = \text{AAS} \times \text{真消化率}$$
\end{keybox}

\section{人体蛋白质营养状况评价}

\begin{defbox}
\textbf{临床检查：}身高、体重、发育；上臂围；上臂肌围、上臂肌区。

\textbf{实验室检查：}
\begin{itemize}[leftmargin=*]
    \item 血清白蛋白（正常值 \imp{35$\sim$50g/L}）
    \item 血清运铁蛋白（正常值 \imp{2.2$\sim$4.0g/L}）
    \item 血清甲状腺结合前蛋白（正常值 \imp{280$\sim$350mg/L}）
    \item 视黄醇结合蛋白（正常值 \imp{26$\sim$76mg/L}）
    \item 尿3-甲基组氨酸：营养不良时排出量减少
    \item 头发的毛干与毛根的形态改变
\end{itemize}

\textbf{蛋白质-能量营养不良（PEM）：}水肿型（Kwashiorkor）/ 干瘦型（Marasmus）。
\end{defbox}

\section{复习题}
\begin{enumerate}[leftmargin=*, itemsep=8pt]
    \item \textbf{【名词解释】}必需氨基酸；条件必需氨基酸；限制氨基酸；氨基酸模式；氮平衡
    \item \textbf{【名词解释】}生物价(BV)；NPU；PER；AAS；PDCAAS
    \item \textbf{【简答题】}简述氮平衡的三种类型并举例。
    \item \textbf{【简答题】}试述蛋白质互补三原则并举谷豆互补例。
    \item \textbf{【简答题】}简述蛋白质营养价值评价的六项指标及公式。
\end{enumerate}

\subsection*{参考答案要点}
\begin{enumerate}[leftmargin=*, itemsep=5pt]
    \item 见正文。EAA 9种（亮/异亮/赖/蛋/苯丙/苏/色/缬/组，组氨酸为婴儿必需）。
    \item 见正文公式。BV=储留氮/吸收氮×100；NPU=消化率×BV=储留氮/食物氮；PER=动物增重/摄入蛋白；AAS=待测AA/理想模式AA×100；PDCAAS=AAS×真消化率。
    \item B=I-(U+F+S)。零（正常成人）、负（老年/饥饿/消耗病）、正（发育/孕期/恢复/增肌）。
    \item 三原则：种属愈远愈好、种类愈多愈好、时间愈近愈好。谷类缺赖氨酸富蛋氨酸，豆类缺蛋氨酸富赖氨酸，同时食用互补。
    \item 含量（氮×6.25）、真消化率/表观消化率、BV、NPU、PER（雄性大鼠28d）、AAS/PDCAAS。
\end{enumerate}
\newpage
